\documentclass[11pt,letterpaper]{article}
\usepackage[T1]{fontenc}
\usepackage[utf8]{inputenc}
\usepackage{lmodern,microtype}
\usepackage[margin=1in,headheight=14pt]{geometry}
\usepackage{amsmath,amssymb,amsthm,mathtools}
\usepackage{booktabs,array,longtable,tabularx}
\usepackage{enumitem}
\usepackage[dvipsnames]{xcolor}
\usepackage{tikz}
\usetikzlibrary{calc,angles,quotes,arrows.meta}
\usepackage{listings}
\usepackage{fancyhdr}
\usepackage{hyperref}
\hypersetup{
  colorlinks=true,
  linkcolor=MidnightBlue,
  urlcolor=BrickRed,
  citecolor=MidnightBlue,
  pdftitle={Mollweide, tangent and power-of-a-point theorems in intuitionistic primitive recursive arithmetic},
  pdfauthor={Prepared for Vladimir Reshetnikov}
}
\pagestyle{fancy}
\fancyhf{}
\fancyhead[L]{\small Constructive triangle and circle identities in iPRA}
\fancyhead[R]{\small 15 November 2026}
\fancyfoot[C]{\thepage}
\setlength{\parskip}{4pt}
\setlength{\parindent}{1.15em}
\setlength{\emergencystretch}{3em}
\setlist{itemsep=3pt,topsep=5pt}
\lstset{basicstyle=\small\ttfamily,breaklines=true,columns=fullflexible,
  frame=single,rulecolor=\color{black!20},showstringspaces=false}

\newtheorem{theorem}{Theorem}[section]
\newtheorem{lemma}[theorem]{Lemma}
\newtheorem{proposition}[theorem]{Proposition}
\newtheorem{corollary}[theorem]{Corollary}
\theoremstyle{definition}
\newtheorem{definition}[theorem]{Definition}
\newtheorem{construction}[theorem]{Construction}
\theoremstyle{remark}
\newtheorem{remark}[theorem]{Remark}
\newtheorem{status}[theorem]{Formalization status}

\newcommand{\ipra}{\mathsf{iPRA}}
\newcommand{\N}{\mathbb N}
\newcommand{\Z}{\mathbb Z}
\newcommand{\Q}{\mathbb Q}
\newcommand{\PR}{\mathrm{PR}}
\newcommand{\Pow}{\operatorname{Pow}}
\newcommand{\Reg}{\operatorname{Reg}}
\newcommand{\Apart}{\operatorname{Apart}}
\newcommand{\norm}[1]{\left\lVert #1\right\rVert}
\newcommand{\abs}[1]{\left|#1\right|}
\newcommand{\ip}[2]{\left\langle #1,#2\right\rangle}
\newcommand{\detm}[2]{\det\!\left(#1,#2\right)}
\newcommand{\eps}{\varepsilon}
\newcommand{\code}[1]{\texttt{#1}}
\newcommand{\half}{\tfrac12}
\newcommand{\sq}[1]{#1^{2}}
\newcommand{\cA}{c_A}
\newcommand{\sA}{s_A}
\newcommand{\cB}{c_B}
\newcommand{\sB}{s_B}
\newcommand{\cC}{c_C}
\newcommand{\sC}{s_C}
\newcommand{\cD}{c_D}
\newcommand{\sD}{s_D}

\title{\textbf{Mollweide's formula, the tangent laws,\\
and power of a point in intuitionistic PRA}\\[6pt]
\large Fully constructive algebraic proofs, exact polynomial certificates,\\
and a Lean completion blueprint}
\author{Prepared for Vladimir Reshetnikov}
\date{15 November 2026}

\begin{document}
\maketitle

\begin{abstract}
We prove seven classical results in a form suitable for intuitionistic
primitive recursive arithmetic: the two Mollweide formulas, the law of
tangents, the law of cotangents, the tangent--secant theorem, the intersecting
chords theorem, the intersecting secants theorem, and their common power-of-a-
point theorem.  The ordinary quotient formulas are not taken as primitive.
Their constructive content is first isolated as denominator-free polynomial
equalities.  Division is introduced only after a positive or apartness witness
has been supplied.

The triangle part uses no analytic theory of radians.  If
\(x=s-a\), \(y=s-b\), and \(z=s-c\), positive square-root witnesses define the
half-angle pairs directly.  Four elementary product equations prove
\(A/2+B/2+C/2=\pi/2\) algebraically; two product-to-sum identities then give
both Mollweide equations and the tangent law.  The inradius is the explicitly
constructed number \(r=K/s\), and the cotangent law reduces to
\(r\cos(A/2)=(s-a)\sin(A/2)\).

The circle part has a single polynomial core.  For a line
\(P+tV\), its power is a quadratic
\(p+2ht+qt^2\).  If two apart parameters \(u,w\) are roots, then
\((w-u)(p-quw)=0\), so constructive cancellation gives \(p=quw\).  All three
classical segment-product theorems are sign/order corollaries.  Tangency is the
orthogonality specialization.

We state precisely what these claims mean inside first-order \(\ipra\): a
uniform finite theorem for natural-number codes, and a schema for every fixed
tuple of primitive-recursive regular-name programs.  A 30-item exact SymPy
check accompanies the paper, but is explicitly separated from formal proof.
The supplied Lean files contain source-level ambient ring certificates, exact
embedded-theory targets for ProveIt's \texttt{ConciseIPRA}, and a module-by-
module route to complete internal formalization.  Lean was unavailable in the
delivery environment, so no Lean compilation is claimed.
\end{abstract}

\section*{Results at a glance}
For a nondegenerate triangle with sides \(a,b,c\), opposite angles
\(A,B,C\), semiperimeter \(s\), and inradius \(r\), we derive
\begin{align*}
 \frac{a+b}{c}
   &=\frac{\cos\frac{A-B}{2}}{\sin\frac C2},
&\frac{a-b}{c}
   &=\frac{\sin\frac{A-B}{2}}{\cos\frac C2},\\
 \frac{a-b}{a+b}
   &=\frac{\tan\frac{A-B}{2}}{\tan\frac{A+B}{2}},
&\frac{\cot\frac A2}{s-a}
   &=\frac{\cot\frac B2}{s-b}
    =\frac{\cot\frac C2}{s-c}=\frac1r,\\
 r^2&=\frac{(s-a)(s-b)(s-c)}s.
\end{align*}
For a circle with center \(O\) and squared radius \(\rho\),
\[
  \Pow(P)=\norm{P-O}^2-\rho
\]
is independent of the chosen secant.  It yields the three familiar forms
\[
 |PT|^2=|PA|\,|PB|,\qquad
 |PA|\,|PD|=|PB|\,|PC|,\qquad
 |AS|\,|SC|=|BS|\,|SD|.
\]
The report proves stronger signed and denominator-free statements first.

\newpage
\begingroup
\small
\setlength{\parskip}{0pt}
\tableofcontents
\endgroup
\newpage

\section{The constructive claim and its exact scope}
\label{sec:scope}

\subsection{Why the usual textbook wording is not yet an iPRA sentence}

First-order \(\ipra\) has variables over natural numbers.  It does not have a
primitive sort containing every real number, every angle, every line, or every
circle.  Nor does it quantify over arbitrary functions \(\N\to\N\).  Thus the
sentence ``for every real triangle'' must be replaced by an arithmetical
statement about numerical codes and fixed primitive-recursive programs.

There are two levels in this report.

\begin{enumerate}
\item \textbf{Uniform finite level.}  Integers, positive-denominator rationals,
finite tuples, sparse polynomials, and finite certificates are coded by natural
numbers.  Every construction is primitive recursive, and every displayed
polynomial equality becomes a decidable equality of coded integers after
clearing positive denominators.
\item \textbf{Regular-name level.}  A particular represented scalar is supplied
by a fixed primitive-recursive program producing rational approximations with a
fixed modulus.  The theorem is a first-order schema: for each fixed tuple of
such programs, \(\ipra\) proves the corresponding universally quantified
formula over their ordinary numerical parameters and requested precision.
\end{enumerate}

This is not a weakening invented for the present problem.  It is the standard
way to express constructive real computation in a first-order arithmetic.  It
also matches the organization of ProveIt's existing triangle-formulas packet
and its \texttt{ConciseIPRA} development.

\subsection{The arithmetic theory}

We use intuitionistic first-order logic with equality, defining equations for
primitive-recursive functions, separation and injectivity of successor, and
induction for quantifier-free formulas with numerical parameters.  Pairing,
finite sequences, bounded search, signed integers, and rational arithmetic are
definitional extensions by fixed primitive-recursive functions.

The current \texttt{ConciseIPRA} repository presents a much smaller native
calculus and proves it equivalent to ordinary \(\ipra\).  For readability, the
human proof is written in ordinary iPRA language.  The final Lean step is to
construct proofs in the repository's \texttt{PRAAx} presentation and transport
them through the existing \texttt{pra\_iff\_native} theorem.  Transport does
not replace the missing proof construction; it is only the last stage.

\subsection{Coded rationals and legitimate finite case distinctions}

A signed integer can be coded, for example, by
\[
  z\longmapsto
  \begin{cases}
    2z,&z\ge0,\\
    -2z-1,&z<0.
  \end{cases}
\]
A rational is a pair \((n,d)\) with signed numerator \(n\) and actual positive
denominator \(d+1\).  Reduction to lowest terms is unnecessary.  Equality and
order are decided by cross multiplication, because both denominators are
positive.  All loops used below are bounded by an explicitly computed natural
number.  A proof by cases on a rational comparison is therefore a proof by
cases on a primitive-recursive Boolean, not excluded middle for an arbitrary
formula.

\subsection{Regular names}

\begin{definition}[Regular dyadic name]
A regular name is a fixed primitive-recursive function \(u:\N\to\Q\) satisfying
\[
  |u(n+1)-u(n)|\le 2^{-n}
\]
for all \(n\).  Two names are extensionally equal when, for every \(k\), their
sufficiently accurate approximations differ by at most \(2^{-k}\).  Any of the
usual equivalent fast-Cauchy conventions can be used; only the constants in the
precision schedules change.
\end{definition}

Addition, subtraction, multiplication on bounded inputs, dot products,
determinants, and fixed polynomials have primitive-recursive realizers with
explicit moduli.  Bounds are stored in the input certificate.  No appeal to an
abstract completeness theorem is needed.

\subsection{Apartness, positivity, and division}

The constructive replacement for ``nonzero'' is positive information.

\begin{definition}[Apartness witness]
A represented scalar \(u\) is apart from zero when the data include a natural
number \(d\) and a proof that \(|u|\ge2^{-d}\).  It is positive when the data
include \(u\ge2^{-d}\).
\end{definition}

An apartness witness yields a primitive-recursive reciprocal algorithm and a
modulus for that algorithm.  We never infer such a witness from the negative
statement \(\neg(u=0)\).  Every cancellation in this paper names the supplied
apart factor.  This is the main logical difference between the constructive
proof and a casual field calculation.

\subsection{Positive square roots by bounded search}

The only analytic-looking operation needed for the triangle identities is the
positive square root.  It has a finitary realization.  For
\(q=p/r\ge0\), \(r\ge1\), and precision \(k\), let \(S=2^k\), choose a simple
bound \(J=(p+1)S\), and scan for the largest \(j\le J\) satisfying
\[
  rj^2\le pS^2.
\]
Set \(\rho_k(q)=j/S\).  The bounded scan is primitive recursive and satisfies
\[
 \rho_k(q)^2\le q<\bigl(\rho_k(q)+2^{-k}\bigr)^2.
 \tag{1.1}\label{eq:root-bracket}
\]
The proof maintains the quantifier-free invariant that \(j\) is the last
successful index.  Applied with an explicit precision schedule to a regular
nonnegative name, this constructs a regular positive-root name.

We use the following constructive uniqueness fact repeatedly.

\begin{lemma}[Positive-root uniqueness]
\label{lem:positive-root}
In an ordered field, if \(u,v\ge0\) and \(u^2=v^2\), then \(u=v\).  In the
coded-name interpretation, the same statement is proved with approximation
bounds; no trichotomy for arbitrary reals is required.
\end{lemma}

\begin{proof}
The factorization \((u-v)(u+v)=0\) is polynomial.  If \(u+v\) is apart from
zero, cancel it.  If no such witness is available, the square equality and
nonnegativity give the quantitative estimate
\(|u-v|^2\le |u^2-v^2|\), after the usual rational interval refinement.  At
precision \(k\), refine both names until their square intervals have width
small enough to force \(|u-v|<2^{-k}\).  This is a bounded procedure.  In all
applications below the quantities are in fact positive, so the simpler
apartness cancellation branch applies.
\end{proof}

\subsection{The proof discipline}

Every result below is organized into three layers.

\begin{enumerate}
\item A \emph{polynomial core}, valid over a commutative ring and checked by
normalization.
\item An \emph{ordered constructive corollary}, where signs and supplied
apartness witnesses justify square-root uniqueness, cancellation, and division.
\item An \emph{iPRA realization}, where all objects are codes or fixed
primitive-recursive names and every loop uses quantifier-free induction.
\end{enumerate}

This separation makes the logical strength visible.  In particular, the
quotient form is never silently substituted for the polynomial theorem.

\section{Finite analytic geometry without analytic angles}
\label{sec:geometry}

\subsection{Vectors, length squares, and oriented area}

For a scalar ring, write a point as \(P=(p_x,p_y)\), and define
\[
 \ip{U}{V}=u_xv_x+u_yv_y,
 \qquad
 \detm{U}{V}=u_xv_y-u_yv_x,
 \qquad
 \norm U^2=\ip U U.
\]
All three are fixed polynomials.  The identities
\begin{align}
 \norm{U-V}^2&=\norm U^2+\norm V^2-2\ip U V,
 \tag{2.1}\label{eq:cos-poly}\\
 \norm U^2\norm V^2&=\ip U V^2+\detm U V^2
 \tag{2.2}\label{eq:gram}
\end{align}
are obtained by expansion.  In a coded rational plane they are decidable
integer equalities after denominator clearing.

For a positively oriented triangle \(ABC\), put
\[
 U=B-A,
 V=C-A,
 \qquad 2K=\detm U V>0.
\]
The side lengths opposite \(A,B,C\) are \(a,b,c\), so
\[
 a^2=\norm{V-U}^2,
 b^2=\norm V^2,
 c^2=\norm U^2.
\]
The determinant apartness witness is the constructive nondegeneracy datum.
Reversing orientation replaces \(2K\) by its positive absolute-value name.

\subsection{Geometric sine and cosine as unit-ray coordinates}

The angle at \(A\) needs no radian measure.  Its geometric cosine and sine are
specified by
\[
  \cos A=\frac{\ip U V}{bc},
  \qquad
  \sin A=\frac{2K}{bc}.
\]
The Gram identity proves \(\cos^2A+\sin^2A=1\).  The denominators are positive
because the triangle certificate contains positive side witnesses.  This pair
is enough for every formula in the report.

A half-angle is represented similarly by a positive pair
\((c_A,s_A)\), not by a real number \(A/2\).  We write the familiar notation
only as a readable abbreviation:
\[
  c_A=\cos(A/2),
 s_A=\sin(A/2).
\]
The defining equations are
\begin{equation}
\begin{gathered}
 c_A^2+s_A^2=1,\qquad
 c_A^2-s_A^2=\cos A,\qquad
 2s_Ac_A=\sin A,\\
 c_A>0,\qquad s_A>0.
\end{gathered}
\tag{2.3}\label{eq:half-angle-characterization}
\end{equation}
Positive-root uniqueness makes this representation unambiguous.

\section{A constructive half-angle certificate for a triangle}
\label{sec:halfangle}

\subsection{Semiperimeter excesses}

Let \(a,b,c>0\) be the sides of a nondegenerate triangle and
\(s=(a+b+c)/2\).  Constructively, strict triangle inequalities are represented
by positive witnesses
\[
 x=s-a,
 y=s-b,
 z=s-c.
 \tag{3.1}\label{eq:xyz}
\]
Equivalently,
\[
 a=y+z,
 b=z+x,
 c=x+y,
 s=x+y+z.
 \tag{3.2}\label{eq:sides-xyz}
\]
These are exact linear identities.  For a rational-coordinate triangle, the
witnesses can be generated uniformly from coordinate bounds and a determinant
apartness witness.  Indeed Heron's polynomial identity below gives
\(K^2=sxyz\).  Compute a positive rational bound \(B\) with
\(s,x,y,z\le B\).  If \(K\ge\delta>0\), then
\[
 \delta^2\le K^2=sxyz\le B^3x,
 \qquad\text{so}\qquad
 x\ge \frac{\delta^2}{B^3},
\]
with cyclic companions.  This argument uses only multiplication and rational
order; it does not divide by \(y\) or \(z\) before their positivity has
been established.  Thus noncollinearity supplies quantitative strictness
rather than only a negated equality.

\subsection{Heron's identity in the needed form}

The coordinate identities \eqref{eq:cos-poly}--\eqref{eq:gram} give
\[
 16K^2=2a^2b^2+2b^2c^2+2c^2a^2-a^4-b^4-c^4.
 \tag{3.3}\label{eq:heron-symmetric}
\]
A ring expansion gives
\[
 (a+b+c)(-a+b+c)(a-b+c)(a+b-c)
 =2a^2b^2+2b^2c^2+2c^2a^2-a^4-b^4-c^4.
 \tag{3.4}\label{eq:heron-factor}
\]
Substituting \eqref{eq:xyz} into the four factors yields
\[
 16K^2=(2s)(2x)(2y)(2z),
 \qquad\text{hence}\qquad
 K^2=sxyz.
 \tag{3.5}\label{eq:heron-xyz}
\]
Everything through \eqref{eq:heron-xyz} is a polynomial certificate.  No
square root has been extracted.

\subsection{Definition of the half-angle components}

Define the six positive components by their squares:
\begin{align}
 s_A^2&=\frac{yz}{bc}, & c_A^2&=\frac{sx}{bc},
 \tag{3.6a}\label{eq:halfA}\\
 s_B^2&=\frac{zx}{ca}, & c_B^2&=\frac{sy}{ca},
 \tag{3.6b}\label{eq:halfB}\\
 s_C^2&=\frac{xy}{ab}, & c_C^2&=\frac{sz}{ab}.
 \tag{3.6c}\label{eq:halfC}
\end{align}
The positive square-root construction of Section~\ref{sec:scope} produces
regular names for them.  The denominators are positive.

\begin{lemma}[Unit equations]
\label{lem:unit-equations}
Each pair is a positive unit pair:
\[
 s_A^2+c_A^2=s_B^2+c_B^2=s_C^2+c_C^2=1.
\]
\end{lemma}

\begin{proof}
For the first pair, clear the positive denominator \(bc\).  The required
numerator identity is
\[
 yz+sx=(z+x)(x+y)=bc,
\]
which follows by expansion using \(s=x+y+z\).  The other two are cyclic.
\end{proof}

\begin{lemma}[The pairs are geometric half-angles]
\label{lem:geometric-halfangles}
The pair \((c_A,s_A)\) satisfies \eqref{eq:half-angle-characterization} for the
geometric angle at \(A\), and cyclically.
\end{lemma}

\begin{proof}
First, by \eqref{eq:heron-xyz} and \eqref{eq:halfA},
\[
 (bc\,s_Ac_A)^2
 =b^2c^2\frac{yz}{bc}\frac{sx}{bc}
 =sxyz=K^2.
\]
Both sides are positive, so Lemma~\ref{lem:positive-root} gives
\[
 bc\,s_Ac_A=K,
 \qquad 2s_Ac_A=\frac{2K}{bc}=\sin A.
 \tag{3.7}\label{eq:half-sine}
\]
Next,
\[
 c_A^2-s_A^2=\frac{sx-yz}{bc}.
\]
A direct expansion gives
\[
 b^2+c^2-a^2=2(sx-yz),
\]
so the cosine-law polynomial \eqref{eq:cos-poly} gives
\[
 c_A^2-s_A^2=\frac{b^2+c^2-a^2}{2bc}=\cos A.
\]
Together with the unit equation and positivity, this is precisely the positive
half-angle characterization.  The proofs for \(B,C\) are cyclic.
\end{proof}

\subsection{Four positive product equations}

The key step is to establish the complementary relation among the three
half-angles without invoking an analytic angle-addition theorem.

\begin{lemma}[Positive product equations]
\label{lem:four-products}
The following equalities hold:
\begin{align}
 c\,s_Ac_B&=y c_C,
 &c\,c_As_B&=x c_C,
 \tag{3.8a}\label{eq:products-cos}\\
 c\,c_Ac_B&=s s_C,
 &c\,s_As_B&=z s_C.
 \tag{3.8b}\label{eq:products-sin}
\end{align}
\end{lemma}

\begin{proof}
All quantities in \eqref{eq:products-cos}--\eqref{eq:products-sin} are
positive.  It is therefore enough to compare squares.
For the first equation, use \eqref{eq:halfA}, \eqref{eq:halfB}, and
\eqref{eq:halfC}:
\[
 c^2s_A^2c_B^2
 =c^2\frac{yz}{bc}\frac{sy}{ca}
 =\frac{s y^2 z}{ab}
 =y^2c_C^2.
\]
Positive-root uniqueness gives \(c s_Ac_B=y c_C\).  Similarly,
\[
 c^2c_A^2s_B^2=x^2c_C^2,
\quad
 c^2c_A^2c_B^2=s^2s_C^2,
\quad
 c^2s_A^2s_B^2=z^2s_C^2.
\]
Applying positive-root uniqueness to each proves the four equations.
\end{proof}

\begin{theorem}[Algebraic complement formula]
\label{thm:complement}
The half-angle pairs satisfy
\begin{align}
 c_C&=s_Ac_B+c_As_B,
 \tag{3.9a}\label{eq:comp-cos}\\
 s_C&=c_Ac_B-s_As_B.
 \tag{3.9b}\label{eq:comp-sin}
\end{align}
Thus, in the conventional notation,
\((A+B)/2=\pi/2-C/2\) at the level of sine/cosine coordinates.
\end{theorem}

\begin{proof}
Add the two equations in \eqref{eq:products-cos}:
\[
 c(s_Ac_B+c_As_B)=(x+y)c_C=cc_C.
\]
The positive side \(c\) is apart from zero, so cancellation gives
\eqref{eq:comp-cos}.  Subtract the two equations in
\eqref{eq:products-sin}:
\[
 c(c_Ac_B-s_As_B)=(s-z)s_C=cs_C,
\]
because \(s-z=c\).  Cancel \(c\) to obtain \eqref{eq:comp-sin}.
\end{proof}

\subsection{A common sine-law scale}

The proof of Lemma~\ref{lem:geometric-halfangles} and its cyclic companions
show
\[
 bc\,s_Ac_A=ca\,s_Bc_B=ab\,s_Cc_C=K.
 \tag{3.10}\label{eq:area-products}
\]
Let \(d>0\) be the constructively defined circumdiameter scale, specified
without division by
\[
 2Kd=abc.
 \tag{3.11}\label{eq:d-scale}
\]
Since \(K\) is apart from zero, its reciprocal name exists.  Multiplying
\eqref{eq:area-products} by the appropriate factor gives
\begin{align}
 a&=2d s_Ac_A,
 &b&=2d s_Bc_B,
 &c&=2d s_Cc_C.
 \tag{3.12}\label{eq:side-halfangle}
\end{align}
This is the extended sine law expressed in half-angle coordinates.  It is the
only scale relation needed for Mollweide's formulas.

\section{Mollweide's two formulas}
\label{sec:mollweide}

Define the signed half-difference pair
\begin{align}
 s_D&=s_Ac_B-c_As_B,
 \tag{4.1a}\label{eq:sD}\\
 c_D&=c_Ac_B+s_As_B.
 \tag{4.1b}\label{eq:cD}
\end{align}
In conventional notation, \(D=(A-B)/2\), so \(s_D=\sin D\) and
\(c_D=\cos D\).  This is not an appeal to an analytic trigonometric library:
\eqref{eq:sD}--\eqref{eq:cD} are the definitions of the pair.  Ring expansion
and the unit equations show \(s_D^2+c_D^2=1\).  Moreover \(c_D>0\), because it
is a sum of positive products.

The following two ring identities are the entire algebraic engine:
\begin{align}
 (s_Ac_B+c_As_B)(c_Ac_B+s_As_B)
 &=s_Ac_A+s_Bc_B,
 \tag{4.2}\label{eq:prod-sum}\\
 (c_Ac_B-s_As_B)(s_Ac_B-c_As_B)
 &=s_Ac_A-s_Bc_B.
 \tag{4.3}\label{eq:prod-diff}
\end{align}
For example, expanding the left side of \eqref{eq:prod-sum} and grouping gives
\[
 s_Ac_A(c_B^2+s_B^2)+s_Bc_B(s_A^2+c_A^2),
\]
which is the right side.  Equation \eqref{eq:prod-diff} is identical with the
middle signs changed.

\begin{theorem}[Denominator-free Mollweide formulas]
\label{thm:mollweide-cross}
For every certified nondegenerate triangle,
\begin{align}
 (a+b)s_C&=c c_D,
 \tag{4.4}\label{eq:moll-sum-cross}\\
 (a-b)c_C&=c s_D.
 \tag{4.5}\label{eq:moll-diff-cross}
\end{align}
\end{theorem}

\begin{proof}
Using \eqref{eq:side-halfangle}, \eqref{eq:prod-sum}, and
\eqref{eq:comp-cos},
\begin{align*}
 a+b
 &=2d(s_Ac_A+s_Bc_B)\\
 &=2d(s_Ac_B+c_As_B)(c_Ac_B+s_As_B)\\
 &=2d c_Cc_D.
\end{align*}
Multiplying by \(s_C\), and using \(c=2d s_Cc_C\), gives
\((a+b)s_C=cc_D\).

Similarly, by \eqref{eq:prod-diff} and \eqref{eq:comp-sin},
\begin{align*}
 a-b
 &=2d(s_Ac_A-s_Bc_B)\\
 &=2d(c_Ac_B-s_As_B)(s_Ac_B-c_As_B)\\
 &=2d s_Cs_D.
\end{align*}
Multiplying by \(c_C\) and again using \(c=2d s_Cc_C\) gives
\((a-b)c_C=cs_D\).  No division has occurred.
\end{proof}

\begin{corollary}[Mollweide's formulas]
\label{cor:mollweide}
With the explicit positive denominator witnesses supplied by the triangle
certificate,
\[
 \frac{a+b}{c}=\frac{c_D}{s_C}
 =\frac{\cos\frac{A-B}{2}}{\sin\frac C2},
 \qquad
 \frac{a-b}{c}=\frac{s_D}{c_C}
 =\frac{\sin\frac{A-B}{2}}{\cos\frac C2}.
 \tag{4.6}\label{eq:mollweide-quotient}
\]
\end{corollary}

\begin{proof}
The quantities \(c,s_C,c_C\) are positive and hence apart from zero.  Divide
\eqref{eq:moll-sum-cross} by \(cs_C\), and divide
\eqref{eq:moll-diff-cross} by \(cc_C\).  Constructively, each division is the
application of the reciprocal program attached to the recorded apartness
witness.
\end{proof}

\begin{remark}
The denominator-free theorem remains meaningful at boundary configurations
where a quotient notation would fail.  The quoted textbook theorem is the
nondegenerate positive corollary, not the primitive statement.
\end{remark}

\section{The law of tangents}
\label{sec:tangents}

The law of tangents is usually obtained by dividing the two Mollweide
formulas.  We retain the cross-multiplied statement as the primary theorem.

\begin{theorem}[Denominator-free law of tangents]
\label{thm:tangent-cross}
With \(s_D,c_D\) as in \eqref{eq:sD}--\eqref{eq:cD},
\[
 (a-b)c_Dc_C=(a+b)s_Ds_C.
 \tag{5.1}\label{eq:tangent-cross}
\]
\end{theorem}

\begin{proof}
Multiply \eqref{eq:moll-diff-cross} by \(c_D\) and
\eqref{eq:moll-sum-cross} by \(s_D\).  Both right sides become
\(cs_Dc_D\), so the left sides are equal.  This is a ring derivation from the
two Mollweide certificates.
\end{proof}

\begin{corollary}[Law of tangents]
\label{cor:tangent-law}
For a nondegenerate triangle,
\[
 \frac{a-b}{a+b}
 =\frac{\tan\frac{A-B}{2}}{\tan\frac{A+B}{2}}.
 \tag{5.2}\label{eq:tangent-law}
\]
\end{corollary}

\begin{proof}
The complement equations \eqref{eq:comp-cos}--\eqref{eq:comp-sin} say that
\[
 \tan\frac{A+B}{2}=\frac{c_C}{s_C}.
\]
The half-difference definition gives
\(\tan((A-B)/2)=s_D/c_D\).  Therefore the right side of
\eqref{eq:tangent-law} is
\[
 \frac{s_D/c_D}{c_C/s_C}=\frac{s_Ds_C}{c_Dc_C}.
\]
The denominator factors \(a+b,c_D,c_C\) are positive.  Dividing
\eqref{eq:tangent-cross} by their product yields the result.  Notice that
\(s_D\) may be negative; no positivity of the numerator is needed.
\end{proof}

\section{The incenter and the law of cotangents}
\label{sec:cotangents}

\subsection{Constructing the incenter}

Let \(A,B,C\) be a positively oriented coordinate triangle and let
\(a,b,c\) be the opposite side lengths.  Define
\[
 I=\frac{aA+bB+cC}{a+b+c},
 \qquad
 r=\frac Ks,
 \tag{6.1}\label{eq:incenter}
\]
or, denominator-free,
\[
 (a+b+c)I=aA+bB+cC,
 \qquad rs=K.
 \tag{6.2}\label{eq:incenter-cross}
\]
The denominator \(a+b+c=2s\) is positive.

\begin{lemma}[Equal signed distances]
\label{lem:incenter-distance}
The point \(I\) has perpendicular distance \(r\) from each side line.  Hence it
is the incenter and \(r\) is the inradius.
\end{lemma}

\begin{proof}
Consider the side \(BC\).  Its length is \(a\).  Determinant linearity and
\eqref{eq:incenter-cross} give
\begin{align*}
 (a+b+c)\detm{C-B}{I-B}
 &=\detm{C-B}{a(A-B)+c(C-B)}\\
 &=a\detm{C-B}{A-B}\\
 &=2aK,
\end{align*}
with the sign chosen by the orientation.  Divide by the positive
\((a+b+c)a=2sa\): the signed distance is \(K/s=r\).  The computations for
\(CA\) and \(AB\) are cyclic.  All divisions are by stored positive side and
perimeter witnesses.  Equivalently, the same conclusion can be recorded with
all denominators cleared.
\end{proof}

\subsection{The basic cotangent equation}

\begin{lemma}[Half-angle/inradius equation]
\label{lem:cot-basic}
With \(x=s-a\), \(y=s-b\), \(z=s-c\),
\[
 r c_A=x s_A,
 \qquad
 r c_B=y s_B,
 \qquad
 r c_C=z s_C.
 \tag{6.3}\label{eq:cot-basic}
\]
\end{lemma}

\begin{proof}
We prove the first equation by comparing positive squares.  From
\(r=K/s\), \eqref{eq:heron-xyz}, and \eqref{eq:halfA},
\begin{align*}
 r^2c_A^2
 &=\frac{K^2}{s^2}\frac{sx}{bc}
 =\frac{sxyz}{s^2}\frac{sx}{bc}
 =x^2\frac{yz}{bc}
 =x^2s_A^2.
\end{align*}
Both \(rc_A\) and \(xs_A\) are positive, so positive-root uniqueness gives
\(rc_A=xs_A\).  The other two equations are cyclic.
\end{proof}

The same equation has the familiar right-triangle interpretation.  Drop a
perpendicular from the incenter to either side adjacent to \(A\).  Its length
is \(r\), while the tangent segment from \(A\) has length \(x=s-a\).  The
angle at \(A\) is bisected, so the right triangle gives
\(\tan(A/2)=r/x\).  Lemma~\ref{lem:cot-basic} is the denominator-free algebraic
form of exactly that statement, and therefore avoids depending on the tangent-
length construction.

\begin{theorem}[Law of cotangents]
\label{thm:cotangent-law}
For a certified nondegenerate triangle,
\[
 \frac{\cot(A/2)}{s-a}
 =\frac{\cot(B/2)}{s-b}
 =\frac{\cot(C/2)}{s-c}
 =\frac1r.
 \tag{6.4}\label{eq:cotangent-law}
\]
\end{theorem}

\begin{proof}
From the first equation of \eqref{eq:cot-basic},
\[
 r c_A=x s_A.
\]
Every factor \(r,x,s_A\) is positive and apart from zero.  Divide by
\(rxs_A\):
\[
 \frac{c_A/s_A}{x}=\frac1r.
\]
Since \(x=s-a\) and \(c_A/s_A=\cot(A/2)\), this is the first equality in
\eqref{eq:cotangent-law}.  The cyclic equations give the other two.
\end{proof}

\begin{corollary}[Inradius formula]
\label{cor:inradius}
\[
 r^2=\frac{(s-a)(s-b)(s-c)}s.
 \tag{6.5}\label{eq:inradius-formula}
\]
\end{corollary}

\begin{proof}
The denominator-free equations are \(rs=K\) and \(K^2=sxyz\).  Squaring the
first gives \(r^2s^2=K^2=sxyz\).  The semiperimeter \(s\) is positive, so
cancel it:
\[
 r^2s=xyz.
\]
Substitute \(x=s-a\), \(y=s-b\), \(z=s-c\), and divide by the positive \(s\).
\end{proof}

\begin{remark}[An independent polynomial certificate]
The accompanying checker also verifies a division-free certificate deriving
\((r^2s-xyz)s_As_Bs_C=0\) from the three equations
\(rc_A=xs_A\), \(rc_B=ys_B\), \(rc_C=zs_C\), the equation
\(s=x+y+z\), and the expanded relation
\(\cos((A+B+C)/2)=0\).  This redundant route is useful for formalization
because it tests the half-angle algebra independently of Heron's shortcut.
\end{remark}

\section{Power of a point: one polynomial theorem}
\label{sec:power}

\subsection{Circle and line charts}

A circle is represented by a center \(O=(o_x,o_y)\) and a nonnegative squared
radius \(\rho\).  No square root is needed.  Define
\[
 \Pow(P)=\norm{P-O}^2-\rho.
 \tag{7.1}\label{eq:power-def}
\]
A line through \(P\) is represented by a direction vector \(V\) with
\(q=\norm V^2>0\):
\[
 L(t)=P+tV.
 \tag{7.2}\label{eq:line-chart}
\]
The direction apartness witness rules out the zero vector.  Parameter
products become actual directed-length products after multiplication by
\(q\), since the directed physical coordinate is \(t\sqrt q\).

Put
\[
 p=\Pow(P),
 \qquad
 h=\ip{P-O}{V},
 \qquad
 q=\norm V^2.
 \tag{7.3}\label{eq:phq}
\]

\begin{lemma}[Line-power expansion]
\label{lem:line-power}
For every scalar parameter \(t\),
\[
 \Pow(L(t))=p+2ht+qt^2.
 \tag{7.4}\label{eq:line-power}
\]
\end{lemma}

\begin{proof}
Expand
\[
 \norm{(P-O)+tV}^2
 =\norm{P-O}^2+2t\ip{P-O}{V}+t^2\norm V^2
\]
after subtracting \(\rho\).  This is a fixed polynomial identity over every
commutative ring.
\end{proof}

\subsection{Two distinct intersections}

Suppose the line meets the circle at the parameters \(u,w\):
\[
 f(u)=p+2hu+qu^2=0,
 \qquad
 f(w)=p+2hw+qw^2=0.
 \tag{7.5}\label{eq:two-roots}
\]
The constructive hypothesis that the intersections are distinct is an
apartness witness for \(w-u\).

\begin{lemma}[Secant factor identity]
\label{lem:secant-factor}
\[
 w f(u)-u f(w)=(w-u)(p-quw).
 \tag{7.6}\label{eq:secant-factor}
\]
\end{lemma}

\begin{proof}
Expand the left side.  The terms \(2huw\) cancel, leaving
\[
 wp-up+q(wu^2-uw^2)
 =(w-u)p-quw(w-u).
\]
Factor \(w-u\).  The calculation is valid over a commutative ring.
\end{proof}

\begin{theorem}[Signed power-of-a-point theorem]
\label{thm:signed-power}
If \(u,w\) are apart roots of the line-circle equation, then
\[
 \Pow(P)=q u w.
 \tag{7.7}\label{eq:signed-power}
\]
Consequently the product of the two directed physical distances from \(P\) to
the intersections is independent of the chosen line and equals
\(\Pow(P)\).
\end{theorem}

\begin{proof}
Substitute \eqref{eq:two-roots} into \eqref{eq:secant-factor}:
\[
 0=(w-u)(p-quw).
\]
The supplied root-separation witness says \(w-u\) is apart from zero.  Cancel
that factor to obtain \(p=quw\).  Since a directed physical coordinate along
\(V\) is \(t\sqrt q\), the product of the two directed distances is
\(u\sqrt q\,w\sqrt q=quw=p\).  No choice of a unit direction is required.
\end{proof}

\begin{remark}[Why the asymmetric combination is useful]
The more familiar Vieta argument says that the product of the roots of
\(qt^2+2ht+p\) is \(p/q\).  Formula \eqref{eq:secant-factor} is better in
\(\ipra\): it avoids dividing by \(q\), and it reveals the exact factor
\(w-u\) whose apartness licenses cancellation.
\end{remark}

\subsection{Tangency}

Let \(T\) be a point on the circle and let \(V\) be a tangent direction:
\[
 \norm{T-O}^2=\rho,
 \qquad
 \ip{T-O}{V}=0.
 \tag{7.8}\label{eq:tangent-data}
\]
Write \(P=T-tV\), so the tangent segment from \(P\) to \(T\) has squared
length \(qt^2\).

\begin{lemma}[Tangent power]
\label{lem:tangent-power}
Under \eqref{eq:tangent-data},
\[
 \Pow(P)=qt^2=\norm{P-T}^2.
 \tag{7.9}\label{eq:tangent-power}
\]
\end{lemma}

\begin{proof}
Expand:
\begin{align*}
 \Pow(T-tV)
 &=\norm{(T-O)-tV}^2-\rho\\
 &=\bigl(\norm{T-O}^2-\rho\bigr)
   -2t\ip{T-O}{V}+t^2\norm V^2\\
 &=qt^2.
\end{align*}
This is again a polynomial identity followed by substitution of the circle and
orthogonality equations.
\end{proof}

There is also a constructive uniqueness statement.  Along the tangent line,
at parameter displacement \(\lambda-t\) from \(T\), the circle residual is
\[
 \Pow\bigl(T+(\lambda-t)V\bigr)=q(\lambda-t)^2.
 \tag{7.10}\label{eq:tangent-residual}
\]
Because \(q>0\), the residual vanishes only when \(\lambda=t\).  Thus the
algebraic tangent condition really produces a single intersection, with no
classical discriminant argument.

\section{The three classical segment-product theorems}
\label{sec:specializations}

The ordinary theorems use unsigned lengths.  Their constructive proofs are
short once the signed theorem is available, but the signs must be stated.

\subsection{Intersecting secants}

Let an external point \(P\) lie on two secant lines.  Orient each line from
\(P\) toward its circle intersections.  Let the first line meet the circle at
parameters \(0<u<w\), and the second at \(0<v<z\).  Then
\[
 |PA|=u\sqrt q,
 |PD|=w\sqrt q,
\]
for the first chart, and similarly on the second.  The signed power theorem
gives
\[
 |PA|\,|PD|=q uw=\Pow(P)=q'vz=|PB|\,|PC|.
\]

\begin{theorem}[Intersecting secants theorem]
\label{thm:secants}
For two secants \(AD\) and \(BC\) meeting outside a circle at \(P\),
\[
 |PA|\,|PD|=|PB|\,|PC|.
 \tag{8.1}\label{eq:secants}
\]
\end{theorem}

\begin{proof}
Apply Theorem~\ref{thm:signed-power} to each line.  The order certificate that
both roots on each oriented line are positive converts the signed product into
the ordinary product of lengths.  Both products equal the same
\(\Pow(P)>0\).
\end{proof}

\subsection{Intersecting chords}

Let \(S\) lie inside the circle.  On a line through \(S\), choose the direction
so that the two intersection parameters are \(-u<0<w\), with \(u,w>0\).
Then
\[
 \Pow(S)=q(-u)w=-q uw<0,
\]
while the unsigned product of the two chord pieces is
\[
 |AS|\,|SC|=(u\sqrt q)(w\sqrt q)=quw=-\Pow(S).
\]
The same is true on every other chord through \(S\).

\begin{theorem}[Intersecting chords theorem]
\label{thm:chords}
For chords \(AC\) and \(BD\) meeting inside a circle at \(S\),
\[
 |AS|\,|SC|=|BS|\,|SD|.
 \tag{8.2}\label{eq:chords}
\]
\end{theorem}

\begin{proof}
Apply the signed power theorem to both chord lines.  Their signed products are
\(\Pow(S)\).  The order certificate for an interior point says the two
parameters on each line have opposite signs.  Negating each signed product
therefore gives the corresponding product of positive segment lengths.  Both
are \(-\Pow(S)\).
\end{proof}

\subsection{Tangent--secant theorem}

Let \(P\) be outside the circle, \(PT\) a tangent, and a secant from \(P\)
meet the circle at \(A,B\), ordered away from \(P\).  Lemma
\ref{lem:tangent-power} and Theorem~\ref{thm:signed-power} give
\[
 |PT|^2=\Pow(P)=|PA|\,|PB|.
\]

\begin{theorem}[Tangent--secant theorem]
\label{thm:tangent-secant}
\[
 |PT|^2=|PA|\,|PB|.
 \tag{8.3}\label{eq:tangent-secant}
\]
\end{theorem}

\begin{proof}
The tangent branch identifies the left side with \(\Pow(P)\).  The external
secant branch identifies the right side with the same power.  The only
cancellations are those already justified in Theorem~\ref{thm:signed-power};
the tangent proof itself uses no cancellation.
\end{proof}

\subsection{The unified power-of-a-point theorem}

\begin{theorem}[Power of a point]
\label{thm:power-final}
For a circle \((O,\rho)\) and a certified point \(P\), the number
\[
 \Pow(P)=\norm{P-O}^2-\rho
\]
has the following constructive meanings.
\begin{enumerate}
\item For any line through \(P\) with two apart circle intersections, it is
  the product of the two directed distances.
\item For an external point and any secant, it is the positive product of the
  two ordinary secant lengths.
\item For an interior point and any chord through the point, its negative is
  the product of the two chord-piece lengths.
\item For an external tangent segment, it is the squared tangent length.
\end{enumerate}
All these values are line-independent because they equal the same polynomial
\(\norm{P-O}^2-\rho\).
\end{theorem}

\begin{proof}
Part 1 is Theorem~\ref{thm:signed-power}.  Parts 2 and 3 are its order/sign
corollaries.  Part 4 is Lemma~\ref{lem:tangent-power}.  The statement of
line-independence is literal equality with the definition \eqref{eq:power-def},
not a separate compactness or continuity argument.
\end{proof}

\section{Exact first-order formulations in iPRA}
\label{sec:ipra}

\subsection{Finite certificate records}

A finite triangle certificate contains natural-number codes for:
\begin{itemize}
\item rational bounds and precision data;
\item side names \(a,b,c\), excess names \(x,y,z\), area name \(K\), and scale
  name \(d\);
\item the six positive half-angle component names;
\item equations \eqref{eq:sides-xyz}, \eqref{eq:heron-xyz},
  \eqref{eq:halfA}--\eqref{eq:halfC}, and \eqref{eq:d-scale};
\item explicit lower bounds for every factor later inverted or cancelled.
\end{itemize}
A finite circle certificate contains codes for \(O,\rho,P,V\), root parameters,
point-on-circle equations, order/sign data, and root-separation witnesses.
Tangency additionally stores the orthogonality equation.

The checker never asks whether an arbitrary represented real is zero.  It only
verifies supplied rational inequalities and polynomial equations to requested
precision.

\subsection{A primitive-recursive polynomial checker}

Represent a monomial by an integer coefficient and a fixed-length exponent
vector.  Represent a sparse polynomial by a finite list of monomials.  The
following operations are primitive recursive:
\begin{enumerate}
\item lexicographic comparison of exponent vectors;
\item insertion into a sorted list;
\item merging equal exponent vectors and deleting zero coefficients;
\item addition, negation, multiplication, and fixed natural powers;
\item evaluation at coded integers or rationals after denominator clearing.
\end{enumerate}

The normalizer's soundness is proved by quantifier-free induction on list
length and on the bounded multiplication loops.  Each identity in this report
is then a closed coefficient-list calculation.  The universal theorem says
that whenever the input residuals vanish, the target residual vanishes.  For
identities presented by explicit ideal certificates, the checker verifies
\[
 T=Q_1H_1+\cdots+Q_mH_m
\]
by normalization.  No Gr\"obner-basis algorithm is needed inside \(\ipra\);
SymPy was used only to discover and independently audit some certificates.

\subsection{The exact finite theorem schema}

For each fixed certificate family, there is a primitive-recursive checker
\(C(\bar n)\) whose output is zero exactly when all input equations and the
claimed output equation are satisfied.  The internal theorem has the shape
\[
 \ipra\vdash\forall\bar n\,\bigl(C_{\mathrm{hyp}}(\bar n)=0
 \to C_{\mathrm{concl}}(\bar n)=0\bigr).
 \tag{9.1}\label{eq:finite-schema}
\]
All induction used in proving \eqref{eq:finite-schema} is on quantifier-free
checker invariants.  Universal closure is applied after the checker proof.

The main denominator-free outputs are:
\begin{align*}
 &(a+b)s_C-cc_D=0,\\
 &(a-b)c_C-cs_D=0,\\
 &(a-b)c_Dc_C-(a+b)s_Ds_C=0,\\
 &rc_A-xs_A=rc_B-ys_B=rc_C-zs_C=0,\\
 &\Pow(P+tV)-\bigl(p+2ht+qt^2\bigr)=0,\\
 &w f(u)-u f(w)-(w-u)(p-quw)=0,\\
 &\Pow(T-tV)-qt^2=0
 \quad\text{under the circle and orthogonality residuals.}
\end{align*}
The quotient statements are separate corollaries whose hypotheses include the
relevant apartness codes.

\subsection{Regular-name lifting}

Fix, externally, a finite tuple of primitive-recursive programs for all input
names.  Their code is not quantified over inside \(\ipra\); the programs are
compiled into the language as fixed function symbols or fixed \(\PR\) terms.
The arithmetic theorem then quantifies only over ordinary parameters and
precision.

For every requested output precision \(k\), compute a larger input precision
\(N(k,B,d)\) from the stored magnitude bound \(B\) and apartness exponent
\(d\).  Evaluate all input names at \(N\), run the rational polynomial
construction, and use the polynomial certificate plus the error-propagation
bound to show the output error is at most \(2^{-k}\).  The schedule is
primitive recursive because the degree and number of operations are fixed.

This gives a first-order theorem such as
\[
 \forall\bar u\,\forall k\,
 \bigl(\operatorname{TriangleCert}(\bar u)
 \to \operatorname{Close}_k((a+b)s_C,cc_D)\bigr),
 \tag{9.2}\label{eq:name-schema}
\]
and similarly for every other identity.  Here \(\operatorname{Close}_k\) is a
coded rational inequality, not equality in a hidden second-order real sort.

\subsection{Quantifier-free induction audit}

The required induction instances are finite and explicit.

\begin{longtable}{p{0.26\textwidth}p{0.65\textwidth}}
\toprule
Component & Quantifier-free invariant \\
\midrule
\endhead
Pair/sequence coding & Projection equations and index bounds for a bounded
scan. \\
Rational arithmetic & Positive denominator and cross-product equations. \\
Square-root mesh & Last successful index, its square below the target, and the
next mesh point above it. \\
Polynomial insertion & Sorted exponent list, coefficient sum preservation,
and absence of duplicate adjacent monomials. \\
Polynomial multiplication & Partial convolution equals the product of the
processed prefixes. \\
Regular-name operations & Integer inequalities bounding the numerator error at
each fixed precision. \\
Reciprocal & A rational interval excludes zero by the supplied apartness bound
and shrinks under bounded refinement. \\
Triangle construction & Finite conjunction of the certificate residuals and
stored positivity bounds. \\
Line-circle construction & The quadratic residual at each supplied parameter,
plus its root-separation inequality. \\
\bottomrule
\end{longtable}

No induction predicate quantifies over all later approximations.  A fixed
precision is a numerical parameter, and the invariant is a decidable relation
on finite codes.

\section{Lean formalization architecture}
\label{sec:lean}

\subsection{What is supplied}

The archive contains three Lean source files and a detailed completion plan.

\begin{description}
\item[\code{AlgebraicCore.lean}] Ambient Mathlib ring identities for the
half-angle algebra, both Mollweide crosses, the tangent cross, an explicit
inradius certificate, the line-power expansion, the secant factorization, and
the tangent residual.
\item[\code{ConstructiveCorollaries.lean}] Ambient ordered-field lemmas showing
exactly which nonzero assumptions are needed for quotient statements and
cancellation.
\item[\code{ObjectTargets.lean}] The exact shape of the deeply embedded
\texttt{PRAAx} targets and the final transport through
\texttt{pra\_iff\_native}.
\item[\code{IPRAFormalizationPlan.md}] A module-by-module route from natural-
number codes through a primitive-recursive polynomial normalizer, regular
names, triangle/circle certificates, object-language proofs, and native export.
\end{description}

The first two files deliberately separate ambient algebra from internal
arithmetic.  A successful \texttt{ring} tactic invocation is evidence for a
polynomial identity, but it is not itself a derivation in \(\ipra\).

\subsection{Core ambient definitions}

A minimal Mathlib-oriented layer needs only scalar coordinates:
\begin{lstlisting}[language={}]
def normSq (x y : R) : R := x^2 + y^2
def dot (x y u v : R) : R := x*u + y*v

def halfDiffSin (sA cA sB cB : R) : R := sA*cB - cA*sB
def halfDiffCos (sA cA sB cB : R) : R := cA*cB + sA*sB
\end{lstlisting}
The polynomial theorems are stated over \code{[CommRing R]}.  Quotient and
cancellation corollaries are stated over \code{[LinearOrderedField R]} with
explicit hypotheses such as \code{c != 0}, \code{sC != 0}, and
\code{w != u}.  This is a useful audit: every field operation in the textbook
form is paired with its constructive premise.

\subsection{Embedded target shape}

The supplied \code{ObjectTargets.lean} follows the API pattern already used in
ProveIt's triangle packet:
\begin{lstlisting}[language={}]
def checkerFormula {arity : Nat} (checkCode : PR arity) : Fm PTm arity :=
  .eq (.app checkCode PTm.var) PTm.zero

def closeAll : (arity : Nat) -> Fm PTm arity -> Fm PTm 0
  | 0, phi => phi
  | arity + 1, phi => closeAll arity (.all phi)

def InternalTarget {arity : Nat} (checkCode : PR arity) : Prop :=
  Prf PRAAx [] (closeAll arity (checkerFormula checkCode))
\end{lstlisting}
The completion obligation is to build the actual \code{PR} syntax trees for
the checkers and construct proofs of these targets.  After that,
\begin{lstlisting}[language={}]
(pra_iff_native (Gamma := []) (phi := phi)).mp h
\end{lstlisting}
transports the theorem to the native concise calculus.

\subsection{Clear path to completion}

A complete formalization can proceed in the following order.

\begin{enumerate}
\item Reuse the existing signed-integer, rational, pairing, and primitive-
recursive infrastructure from \texttt{ConciseIPRA} and the delivered triangle
packet.
\item Implement sparse polynomial syntax and a normalizer as concrete
\code{PR} terms.  Prove soundness in \code{PRAAx} by quantifier-free induction.
\item Compile the fixed identities in \code{AlgebraicCore.lean} into explicit
coefficient-list certificates.  The supplied Python report gives a regression
oracle, not the formal proof.
\item Reuse or complete the bounded positive-square-root implementation from
the triangle packet, and implement reciprocal from apartness.
\item Define \code{TriangleCert}; prove the four positive product equations
and then the Mollweide, tangent, and cotangent targets.
\item Define \code{CircleCert} and \code{LineChart}; prove the quadratic
expansion, factor identity, signed product theorem, tangent theorem, and the
three sign/order specializations.
\item Lift finite rational statements to fixed regular-name programs using
explicit precision schedules.
\item Construct the closed \code{Prf PRAAx} terms and transport them to
\code{NPrf}.
\end{enumerate}

No step requires a new logical principle.  The remaining work is substantial
proof engineering, but each obligation is finite and named.

\begin{status}
Neither \code{lean} nor \code{lake} was installed in the delivery environment.
The Lean sources were therefore not compiled.  They contain no
\code{sorry}, but syntax/type correctness has not been mechanically confirmed.
No completed internal \code{Prf PRAAx} proof object is claimed.
\end{status}

\section{Executed exact checks and worked examples}
\label{sec:verification}

\subsection{Symbolic audit}

The script \code{verification/check\_identities.py} performs 30 exact checks.
It uses SymPy polynomial expansion, rational cancellation, and, for two
Mollweide consequences, ideal reduction.  All numerical examples use Python's
\code{fractions.Fraction}; no floating-point arithmetic occurs.  The captured
result is
\begin{quote}
\code{ALL CERTIFICATES PASS (30 checks)}.
\end{quote}
The checks include:
\begin{itemize}
\item the semiperimeter substitution and all three half-angle unit numerators;
\item the four positive-product square identities;
\item the three area/half-angle products;
\item both Mollweide crosses and the tangent cross;
\item the cotangent right-triangle equations and an independent inradius
certificate;
\item the line-power expansion, secant factor identity, tangent decomposition,
and tangent uniqueness residual;
\item exact rational examples of tangent--secant and intersecting chords.
\end{itemize}

These computations are a development audit.  They are not substituted for the
human proof or for the proposed \(\ipra\) proof objects.

\subsection{Rational tangent--secant example}

Take the circle \(x^2+y^2=25\) and the external point \(P=(13,0)\).  The point
\[
 T=\left(\frac{25}{13},\frac{60}{13}\right)
\]
lies on the circle.  The radius \(OT\) is perpendicular to \(PT\), and exact
fraction arithmetic gives
\[
 |PT|^2=144=\Pow(P)=13^2-25.
\]
A rational secant through
\[
 A=(3,4),
 \qquad
 B=\left(\frac{17}{29},\frac{144}{29}\right)
\]
has chart direction \(V=A-P=(-10,4)\), so \(q=116\).  Its two root parameters
are \(1\) and \(36/29\), whence
\[
 q\cdot1\cdot\frac{36}{29}=116\cdot\frac{36}{29}=144.
\]
Thus the tangent square and the secant product agree exactly.

\subsection{Rational intersecting-chords example}

On the same circle let \(S=(1,0)\).  One chord is the diameter along the
\(x\)-axis, with pieces \(6\) and \(4\), product \(24\).  A second line passes
through
\[
 C_1=(3,4),
 \qquad
 C_2=\left(-\frac75,-\frac{24}{5}\right).
\]
For the chart direction \(V=C_1-S=(2,4)\), \(q=20\), and the root-parameter
product is \(-6/5\).  Hence
\[
 \Pow(S)=20\left(-\frac65\right)=-24.
\]
The unsigned product of the two pieces is \(-\Pow(S)=24\), equal to the
diameter-chord product.

\section{Logical audit and limitations}
\label{sec:audit}

\subsection{Principles not used}

The proofs do not use:
\begin{itemize}
\item excluded middle for arbitrary formulas;
\item choice or countable choice;
\item a completed real-number object quantified over as a second-order domain;
\item an unrestricted total evaluator for arbitrary primitive-recursive program
codes;
\item compactness, completeness of the reals, or an analytic theory of radians;
\item division by a merely nonzero quantity without a positive apartness
witness;
\item cancellation of the two-root factor without a root-separation witness.
\end{itemize}

\subsection{What is assumed as constructive input}

A nondegenerate triangle carries quantitative strictness: positive side,
area, and semiperimeter-excess witnesses.  A secant with two distinct
intersections carries root-separation data.  An external/interior
specialization carries the finite order signs that identify the signed product
with the corresponding unsigned product.  These are not additional classical
axioms; they are the positive data hidden by ordinary textbook phrases such as
``nondegenerate'', ``distinct'', ``outside'', and ``inside''.

\subsection{What remains formalization work}

The human mathematical argument is complete for the stated coded/certified
constructive formulations.  The exact symbolic audit was executed.  The PDF
was compiled from the supplied source.  The deeply embedded Lean proof in
\texttt{ConciseIPRA} remains to be implemented: concrete \code{PR} checker
codes, their \code{PRAAx} soundness derivations, regular-name lifting, and the
final native transport.  The archive labels this boundary in every Lean file.

\section{Further research questions}

\begin{enumerate}
\item Can the triangle half-angle certificate be reduced to fewer positive-root
names by using only \((s_A,c_A)\), \((s_B,c_B)\), and deriving the third pair
with a certified normalization step?
\item What is the smallest standalone \texttt{ConciseIPRA} proof checker for
the fixed polynomial certificates in this packet?
\item Can the signed power theorem be generalized in the same checker to
quadratic forms and spheres in arbitrary fixed dimension without increasing the
logical layer?
\item Can root separation for a rational line-circle intersection be generated
uniformly from a nonzero discriminant witness, yielding an especially compact
secant certificate?
\item Is there a proof-friendly projective formulation in which the secant,
chord, and tangent cases become one homogeneous determinant identity while
retaining transparent constructive sign corollaries?
\end{enumerate}

\appendix

\section{Certificate inventory}
\label{app:certificates}

The executed audit records the following checks.  Names match
\code{verification/results.json}.

\begin{longtable}{r p{0.74\textwidth}}
\toprule
No. & Certificate \\
\midrule
\endhead
1--3 & \code{semiperimeter\_x}, \code{semiperimeter\_y},
\code{semiperimeter\_z}. \\
4--6 & The three half-angle unit numerators. \\
7--10 & The four squared positive-product identities of
Lemma~\ref{lem:four-products}. \\
11--13 & The three area/half-angle product identities
\(bc s_Ac_A=K\) and cyclic companions, checked after squaring. \\
14--15 & The two denominator-free Mollweide identities. \\
16--17 & Derivation and direct ideal check of the tangent cross. \\
18--20 & The three cotangent right-triangle cross identities. \\
21 & A full explicit polynomial certificate for the inradius formula. \\
22 & The line-power expansion. \\
23 & The secant factor identity. \\
24 & Tangent power decomposition. \\
25 & Tangent-line uniqueness residual. \\
26--28 & Algebraic cores of intersecting secants, intersecting chords, and
tangent--secant. \\
29--30 & Exact rational tangent--secant and chord examples. \\
\bottomrule
\end{longtable}

\section{A compact theorem interface for formalization}
\label{app:interface}

A Lean or proof-checker implementation can expose the following interface.

\begin{description}
\item[\code{TriangleCert}] Fields \(a,b,c,s,x,y,z,K,d\), six half-angle
components, magnitude bounds, positivity witnesses, and all defining residuals.
\item[\code{mollweideSumCross}] Returns
\((a+b)s_C=cc_D\).
\item[\code{mollweideDiffCross}] Returns
\((a-b)c_C=cs_D\).
\item[\code{tangentCross}] Returns
\((a-b)c_Dc_C=(a+b)s_Ds_C\).
\item[\code{cotangentCrossA/B/C}] Returns
\(rc_A=xs_A\) and cyclic companions.
\item[\code{Circle}] Fields \(O_x,O_y,\rho\).
\item[\code{LineChart}] Fields \(P_x,P_y,V_x,V_y\), with
\(V_x^2+V_y^2\) apart from zero.
\item[\code{powerLine}] Returns \(f(t)=p+2ht+qt^2\).
\item[\code{signedSecantProduct}] Given two root residuals and separation,
returns \(p=quw\).
\item[\code{tangentPower}] Given point-on-circle and orthogonality residuals,
returns \(p=qt^2\).
\item[\code{unsignedSpecializations}] Uses order certificates to return the
three textbook segment-product equations.
\end{description}

The interface places all partiality in certificate construction.  Once a
certificate is present, every theorem is total and proof checking is finite.

\begin{thebibliography}{99}

\bibitem{mollweide-wiki}
Wikipedia contributors,
\emph{Mollweide's formula}, standard statement and relation to the tangent law,
accessed 15 November 2026.
\url{https://en.wikipedia.org/wiki/Mollweide's_formula}

\bibitem{tangent-wiki}
Wikipedia contributors,
\emph{Law of tangents}, standard statement,
accessed 15 November 2026.
\url{https://en.wikipedia.org/wiki/Law_of_tangents}

\bibitem{cotangent-wiki}
Wikipedia contributors,
\emph{Law of cotangents}, standard statement and inradius formula,
accessed 15 November 2026.
\url{https://en.wikipedia.org/wiki/Law_of_cotangents}

\bibitem{tangent-secant-wiki}
Wikipedia contributors,
\emph{Tangent--secant theorem}, standard segment-product statement,
accessed 15 November 2026.
\url{https://en.wikipedia.org/wiki/Tangent--secant_theorem}

\bibitem{chords-wiki}
Wikipedia contributors,
\emph{Intersecting chords theorem}, standard statement,
accessed 15 November 2026.
\url{https://en.wikipedia.org/wiki/Intersecting_chords_theorem}

\bibitem{secants-wiki}
Wikipedia contributors,
\emph{Intersecting secants theorem}, standard statement,
accessed 15 November 2026.
\url{https://en.wikipedia.org/wiki/Intersecting_secants_theorem}

\bibitem{power-wiki}
Wikipedia contributors,
\emph{Power of a point}, definition and common theorem,
accessed 15 November 2026.
\url{https://en.wikipedia.org/wiki/Power_of_a_point}

\bibitem{proveit}
V. Reshetnikov,
\emph{ProveIt}, commit
\code{948fe06760595e3bce468b67d25032bbbf7f376d}, 15 November 2026.
\url{https://github.com/VladimirReshetnikov/ProveIt}

\bibitem{concise-ipra}
V. Reshetnikov,
\emph{Concise presentations of intuitionistic primitive recursive arithmetic},
\code{Logic/PrimitiveRecursiveArithmetic/ConciseIPRA} in \cite{proveit}.

\bibitem{triangle-packet}
V. Reshetnikov,
\emph{Heron, sine, and cosine laws in intuitionistic PRA},
\code{Logic/PrimitiveRecursiveArithmetic/docs/proofs/triangle\_formulas\_ipra}
in \cite{proveit}.

\bibitem{geometry-packet}
V. Reshetnikov,
\emph{Eight classical geometry theorems in iPRA},
\code{Logic/PrimitiveRecursiveArithmetic/docs/proofs/eight\_geometry\_theorems\_ipra}
in \cite{proveit}.

\end{thebibliography}

\end{document}
